\documentclass[a4paper,11pt]{article}
\usepackage[utf8]{inputenc}
\usepackage[T1]{fontenc}
\usepackage[french]{babel}
\usepackage[top=2cm, bottom=2cm, left=2.2cm, right=2.2cm]{geometry}
\usepackage{xcolor}
\usepackage{booktabs}
\usepackage{tabularx}
\usepackage{enumitem}
\usepackage{listings}
\usepackage{float}
\usepackage{amssymb}
\usepackage{pifont}
\usepackage{tcolorbox}
\usepackage{hyperref}
\hypersetup{colorlinks=true, linkcolor=primary, urlcolor=primary}

\definecolor{primary}{HTML}{1A237E}
\definecolor{success}{HTML}{2E7D32}
\definecolor{warning}{HTML}{E65100}
\definecolor{lightgreen}{HTML}{E8F5E9}
\definecolor{lightblue}{HTML}{E3F2FD}
\definecolor{codebg}{HTML}{F5F5F5}
\definecolor{codeblue}{rgb}{0.13,0.29,0.53}
\definecolor{codegreen}{rgb}{0.0,0.50,0.0}
\definecolor{codegray}{rgb}{0.5,0.5,0.5}

\lstdefinestyle{sh}{backgroundcolor=\color{codebg}, basicstyle=\ttfamily\footnotesize,
  keywordstyle=\color{codeblue}\bfseries, commentstyle=\color{codegreen}\itshape,
  numberstyle=\tiny\color{codegray}, breaklines=true, frame=single,
  showstringspaces=false, tabsize=2, columns=fullflexible}
\lstset{style=sh}
\newcommand{\ok}{\textcolor{success}{\checkmark}}

\title{\textbf{Communication V2V par Vanetza (ETSI CAM)\\ entre PC et Jetson --- ProtoVA}\\[2mm]
\large Échange de messages C-ITS sur WiFi + intégration ROS2 et alerte de proximité}
\author{Douae Sebti}
\date{25 juin 2026}

\begin{document}
\maketitle

\begin{abstract}
\noindent Ce document décrit la mise en place d'une communication \textbf{V2V}
(\emph{vehicle-to-vehicle}) entre un PC et une Jetson, au standard européen
\textbf{ETSI C-ITS} (messages \textbf{CAM}), à l'aide de la pile open-source
\textbf{Vanetza}. On montre que de vrais messages CAM sont échangés sur un simple
réseau \textbf{WiFi} (sans matériel radio 802.11p), puis l'intégration dans
\textbf{ROS2} via un node qui publie les véhicules voisins et déclenche une
\textbf{alerte de proximité}.
\end{abstract}

\tableofcontents
\bigskip

\section{Contexte et choix techniques}
Le V2V permet aux véhicules d'échanger leur état (position, vitesse, cap) et des
alertes, pour la sécurité coopérative.
\begin{itemize}[leftmargin=1.4em]
  \item \textbf{Standard.} En France/Europe, c'est l'\textbf{ETSI C-ITS} :
        \textbf{CAM} (\emph{Cooperative Awareness Message}, position/vitesse
        périodiques) et \textbf{DENM} (alertes d'événements) --- et non le
        \textbf{BSM} américain (SAE J2735).
  \item \textbf{Outil : Vanetza} (\url{https://github.com/riebl/vanetza}),
        implémentation open-source de la pile ETSI ITS-G5 (GeoNetworking, BTP,
        CAM/DENM ASN.1, sécurité).
  \item \textbf{Sans radio dédiée.} Le vrai V2V utilise des radios 802.11p /
        C-V2X. On ne les a pas (les modems 5G Quectel sont \emph{Uu}, pas
        \emph{PC5}). L'application \texttt{socktap} de Vanetza tourne sur une
        \textbf{interface réseau normale} (socket UDP/raw) : deux machines sur le
        même \textbf{WiFi} échangent de vrais CAM. C'est correct au niveau
        protocole (messages ETSI), seule la couche radio diffère.
\end{itemize}

\section{Compilation de Vanetza}
Sur le PC et la Jetson (Ubuntu 22.04, ROS2 Humble) :
\begin{lstlisting}[language=bash]
git clone --depth 1 https://github.com/riebl/vanetza.git
cd vanetza && mkdir build && cd build
cmake -DBUILD_SOCKTAP=ON ..        # socktap est OFF par defaut
make -j4                            # backend crypto = OpenSSL (pas besoin de Crypto++)
# binaire : ~/vanetza/build/bin/socktap
\end{lstlisting}
\textbf{Notes :} Vanetza embarque son propre code GeographicLib et utilise
\textbf{OpenSSL} (présent) comme backend crypto --- aucune dépendance
supplémentaire à installer. L'option \texttt{BUILD\_SOCKTAP=ON} est indispensable
(sinon seules les bibliothèques sont compilées, pas l'application).

\section{Échange de CAM (le cœur du V2V)}
On lance \texttt{socktap} sur chaque machine, en lien \textbf{UDP} (pas de
\texttt{sudo}, contrairement au mode \emph{ethernet} raw) :
\begin{lstlisting}[language=bash]
# Vehicule A (PC),     interface wlp2s0
./bin/socktap -l udp -i wlp2s0 -a ca --print-rx-cam --security none --station-id 1
# Vehicule B (Jetson), interface wlP1p1s0
./bin/socktap -l udp -i wlP1p1s0 -a ca --print-rx-cam --security none --station-id 2
\end{lstlisting}
\textbf{Résultat (\ok) : échange bidirectionnel vérifié.}
\begin{center}
\begin{tabularx}{\textwidth}{@{}l l X@{}}
\toprule
Sens & Reçu & Preuve \\
\midrule
Jetson $\rightarrow$ PC & le PC voit \texttt{Station ID: 2} & paquet depuis MAC \texttt{14:75:5b:15:59:32} \\
PC $\rightarrow$ Jetson & la Jetson voit \texttt{Station ID: 1} & paquet depuis MAC \texttt{10:e1:8e:56:07:1a} \\
\bottomrule
\end{tabularx}
\end{center}
Chaque CAM fait $\sim$99 octets, émis à $\sim$1 Hz, au format ETSI (Protocol
Version 2, conteneur \emph{CoopAwareness} : position GPS, cap, vitesse, type et
dimensions du véhicule).

\begin{tcolorbox}[colback=lightgreen,colframe=success,title=Points clés]
Vrais messages \textbf{ETSI CAM} échangés \textbf{sur WiFi}, \textbf{sans 802.11p}
et \textbf{sans sudo} (lien UDP). C'est la communication V2V de base
opérationnelle.
\end{tcolorbox}

\section{Intégration ROS2 : node \texttt{v2v\_node.py}}
Pour exploiter le V2V dans le robot, un node ROS2 (Python, \texttt{rclpy}) fait le
pont : il \textbf{lance socktap}, \textbf{lit ses CAM reçus} (sortie
\texttt{--print-rx-cam}), les publie en topics ROS2 et calcule une \textbf{alerte
de proximité}. Aucune modification C++ de Vanetza n'est nécessaire.

\subsection{Architecture}
\begin{center}
\texttt{socktap (CAM ETSI, WiFi)} $\rightarrow$ \texttt{v2v\_node.py} (parse)
$\rightarrow$ \texttt{/v2v/remote\_vehicles} + \texttt{/v2v/alert}
\end{center}
\begin{itemize}[leftmargin=1.4em]
  \item \textbf{Parsing.} À chaque CAM reçu : \texttt{Station ID},
        \texttt{Longitude}/\texttt{Latitude} (en $10^{-7}$ degré),
        \texttt{Heading} ($0{,}1^\circ$), \texttt{Speed} ($0{,}01$ m/s).
  \item \textbf{Distance.} Formule de \emph{haversine} entre la position du
        véhicule et celle reçue.
  \item \textbf{Topics.} \texttt{/v2v/remote\_vehicles} (JSON : id, lat, lon, cap,
        vitesse, \texttt{distance\_m}) ; \texttt{/v2v/alert} (texte) si un
        véhicule est à moins de \texttt{alert\_dist} mètres.
  \item \textbf{Paramètres.} \texttt{interface, station\_id, my\_lat, my\_lon,
        alert\_dist, socktap\_bin}.
\end{itemize}
\begin{lstlisting}[language=bash]
python3 ~/v2v_node.py --ros-args \
  -p interface:=wlP1p1s0 -p station_id:=2 \
  -p my_lat:=48.7668 -p my_lon:=11.432 -p alert_dist:=10.0 \
  -p socktap_bin:=/home/protova2/vanetza/build/bin/socktap
\end{lstlisting}

\subsection{Résultats (\ok)}
Véhicule B (Jetson) à (48.76680, 11.43200), seuil d'alerte 10 m :
\begin{center}
\begin{tabular}{@{}lll@{}}
\toprule
Cas & \texttt{/v2v/remote\_vehicles} & \texttt{/v2v/alert} \\
\midrule
PC proche ($\sim$5 m) & \texttt{distance\_m: 5.56} & \textbf{ALERTE véhicule 1 à 5.6 m} \\
PC loin ($\sim$1,1 km) & \texttt{distance\_m: 1111.95} & (aucune) \\
\bottomrule
\end{tabular}
\end{center}
La logique discrimine correctement proche/loin. Les topics (petits messages
texte) traversent bien le WiFi.

\section{Position dynamique (\ok)}
Par défaut socktap émet une position \emph{statique}. Pour que le CAM porte la
position \textbf{réelle/évolutive} du véhicule, on a ajouté à socktap un
\textbf{fournisseur de position alimenté par UDP} :
\begin{itemize}[leftmargin=1.4em]
  \item nouveau \texttt{udp\_position\_provider.cpp/.hpp} (option \texttt{-p udp},
        port \texttt{--pos-port} 9001) : il reçoit des datagrammes
        \texttt{"lat,lon"} (degrés) et met à jour la position courante ;
  \item le node \texttt{v2v\_node.py} s'abonne au topic ROS2
        \texttt{/v2v/my\_gps} (\texttt{Float64MultiArray [lat, lon]}) et pousse la
        position à socktap à 5 Hz $\rightarrow$ elle part dans le CAM émis ;
  \item la position peut venir de l'odométrie/GPS réel, ou d'un simulateur
        (\texttt{v2v\_gps\_sim.py}) pour une démo.
\end{itemize}
\textbf{Démo validée :} un véhicule mobile (PC) se rapproche de 100 m à 3,3 m ; le
véhicule fixe (Jetson) suit la distance en temps réel et déclenche l'alerte au
passage sous 10 m. \emph{Important :} les deux véhicules tournent sur des
\textbf{domaines ROS2 distincts} (le lien V2V passe par Vanetza/WiFi, pas par ROS)
pour que leurs topics ne se mélangent pas.

\section{Limites et perspectives}
\begin{itemize}[leftmargin=1.4em]
  \item \textbf{DENM} : ajouter les messages d'événement (alerte de freinage
        d'urgence, obstacle).
  \item \textbf{Réaction physique} : relier l'alerte V2V à l'AEB (freinage auto).
  \item \textbf{Sécurité} : passer en \texttt{--security certs} (messages signés,
        comme le C-ITS réel).
\end{itemize}

\paragraph{Livrables.}
\texttt{\textasciitilde/vanetza/} (build + \texttt{bin/socktap}, PC et Jetson, avec
le provider de position UDP) ; \texttt{\textasciitilde/v2v\_node.py} (node ROS2
pont/alerte, PC et Jetson) ; \texttt{\textasciitilde/v2v\_gps\_sim.py} (simulateur
de déplacement).

\end{document}
